\chapter*{Manuals}
\addcontentsline{toc}{chapter}{Manuals}

\section*{Installation Manual}
\addcontentsline{toc}{section}{Installation Manual}

This manual describes the steps required to install and prepare the local environment used for the RAG experiments.

\subsection*{System Requirements}
The project was executed locally using Python, Jupyter Notebook, Ollama, and Milvus. The experiments require enough memory and storage to load the document corpus, generate embeddings, and execute the RAG pipelines.

\subsection*{Python Environment}
A Python virtual environment should be created before installing the project dependencies. After activating the environment, the required packages can be installed from the project requirements file.

\subsection*{Model Setup}
The experiments use Ollama for local model serving. The generation model used in the thesis is \texttt{llama3.1}, and the embedding model is \texttt{mxbai-embed-large}. Both models must be available before running the notebooks.

\subsection*{Milvus Setup}
The vector-database, tuned-vector-database, hybrid, reranking, and hybrid-reranking experiments require a local Milvus service. The service is started using the project script \texttt{scripts/standalone\_embed.sh}.

\subsection*{Project Data}
The OrchestrAI document corpus must remain in the expected project directory structure so that each notebook can load the same input files.

\newpage
\section*{User Manual}
\addcontentsline{toc}{section}{User Manual}

This manual explains how to run the implemented RAG pipelines and generate the experimental outputs.

\subsection*{Running the Notebooks}
The notebooks should be executed in the following order:
\begin{enumerate}
    \item \texttt{rag\_basic.ipynb}
    \item \texttt{rag\_vectordb.ipynb}
    \item \texttt{rag\_vectordb\_tuned.ipynb}
    \item \texttt{rag\_reranker.ipynb}
    \item \texttt{rag\_hybrid.ipynb}
    \item \texttt{rag\_hybrid\_\&\_rerank.ipynb}
\end{enumerate}

Each notebook loads the same document corpus, executes the same five scenario questions, records retrieved sources and generated answers, and writes timing information to the shared results file.

\subsection*{Evaluation Workflow}
After the notebooks have been executed, the generated results are used to create the manual evaluation sheet. Manual correctness scores are then added, and the deterministic evaluation script calculates retrieval and answer-quality metrics.

\subsection*{Generating Figures}
The visualization notebook reads the processed results and generates figures for answer quality, retrieval quality, timing breakdown, scenario comparison, and radar-chart summaries.

\subsection*{Interpreting Outputs}
The most important outputs are the generated answers, retrieved sources, timing fields, manual correctness scores, retrieval metrics, and automated answer-quality metrics. These outputs are used in the results chapter to compare the six RAG pipeline variants.

